\documentclass[11pt]{scrartcl}

\usepackage[stixtwo]{fontsetup}
\usepackage{amsmath}

\begin{document}
Teljes indukcióval belátjuk, hogy $2 \le a_n$ minden $n$-re.

$n = 1$ esetén $2 \le a_1 = 13$ teljesül.

Ha $2 \le a_n$, akkor $4 \le a_n^2$, $12 \le a_n^2 + 8$ és
$4 \le \frac{a_n^2 + 8}{3}$, tehát
\(
2 \le \sqrt{\frac{a_n^2 + 8}{3}} = a_{n + 1}
\).

Megmutatjuk, hogy $a_n \ge a_{n + 1}$ minden $n$-re.

Már beláttuk, hogy $a_n \ge 2$. Innen $2a_n^2 \ge 8$, $3a_n^2 \ge a_n^2 + 8$,
$a_n^2 \ge \frac{a_n^2 + 8}{3}$. Mivel $a_n$ pozitív, ezért gyökvonás után
kapjuk az állítást.

Mivel $a_n$ monoton csökkenő, így felülről korlátos, és
$\sup a_n = a_1 = 13$.

Megmutattuk, hogy $a_n$ monoton és korlátos, így a konvergencia elégséges
feltétele miatt konvergens.

Legyen $A = \lim a_n \in \mathbb{R}$! Felhasználva a részsorozat
határértékére és a határérték és a műveletek kapcsolatára vonatkozó
tételeket, kapjuk az \(A = \sqrt{\frac{A^2 + 8}{3}}\) egyenletet.

Mivel $a_n$ monoton csökken, ezért $\inf a_n = \lim a_n = 2$.
\end{document}
